\section{Mathematical foundations}

This chapter assumes no prior quantum mechanics. The essential ideas are
numbers, vectors, probability, and transformations. Quantum notation can look
unfamiliar, but each symbol has a concrete computational meaning.

\subsection{Real numbers and complex numbers}

The real numbers include familiar values such as $-2$, $0$, $1/3$, and
$\sqrt2$. A real number can be placed on a one-dimensional number line.
Quantum amplitudes require a second numerical direction, represented by the
\emph{imaginary unit}
\[
i=\sqrt{-1},\qquad i^2=-1.
\]
There is no real number whose square is $-1$, so $i$ extends the number system.
A complex number has the form
\[
z=a+bi,
\]
where $a$ is the real part and $b$ is the imaginary part. Examples are
$3+2i$, $-i$, and $1$ (which is $1+0i$).

Complex arithmetic follows ordinary algebra together with $i^2=-1$:
\[
(a+bi)+(c+di)=(a+c)+(b+d)i,
\]
\[
(a+bi)(c+di)=(ac-bd)+(ad+bc)i.
\]
For example,
\[
(1+i)^2=1+2i+i^2=2i.
\]

\subsection{The complex plane, magnitude, and conjugate}

A complex number is a point or arrow in a plane: horizontal coordinate $a$,
vertical coordinate $b$. Its magnitude is the arrow length,
\[
|z|=\sqrt{a^2+b^2}.
\]
The complex conjugate reflects the arrow across the real axis:
\[
z^*=a-bi.
\]
Multiplying by the conjugate gives a nonnegative real number:
\[
z^*z=(a-bi)(a+bi)=a^2+b^2=|z|^2.
\]
Quantum measurement probabilities use $|z|^2$, not $z$ itself. This guarantees
real, nonnegative probabilities even though amplitudes may be complex.

\subsection{Angles, radians, sine, and cosine}

An arrow in the complex plane may be described by its length $r$ and angle
$\theta$:
\[
a=r\cos\theta,\qquad b=r\sin\theta.
\]
Quantum formulas normally measure angles in radians. One full turn is $2\pi$
radians, a half turn is $\pi$, and a quarter turn is $\pi/2$. Conversion is
\[
\theta_{\rm radians}=\theta_{\rm degrees}\frac{\pi}{180}.
\]
This is why \code{PHASE=90d} and \code{PHASE=pi/2} describe the same angle.

\subsection{Euler's formula and Euler's identity}

Euler's formula connects complex numbers to rotations:
\[
e^{i\theta}=\cos\theta+i\sin\theta.
\]
The value $e^{i\theta}$ always has magnitude one. As $\theta$ changes, it moves
around the unit circle in the complex plane. Important cases are
\[
e^{i0}=1,\quad e^{i\pi/2}=i,\quad e^{i\pi}=-1,\quad
e^{i2\pi}=1.
\]
Setting $\theta=\pi$ gives Euler's identity,
\[
e^{i\pi}+1=0.
\]
In a phase gate, multiplying an amplitude by $e^{i\theta}$ rotates that
amplitude without changing its magnitude. Its immediate measurement
probability is unchanged, but its relationship to other amplitudes changes.
Later gates can convert that relative phase into different probabilities.

\subsection{Vectors and basis vectors}

A column vector is an ordered list of numbers. The two basis vectors used for
one qubit are
\[
\ket0=\begin{bmatrix}1\\0\end{bmatrix},\qquad
\ket1=\begin{bmatrix}0\\1\end{bmatrix}.
\]
The vertical bars and angle bracket form Dirac's \emph{ket} notation.
\code{0p} prepares $\ket0$ and \code{1p} prepares $\ket1$.
Any one-qubit pure state is a complex linear combination:
\[
\ket{\psi}=\alpha\ket0+\beta\ket1
=\begin{bmatrix}\alpha\\\beta\end{bmatrix}.
\]
The complex numbers $\alpha$ and $\beta$ are \emph{probability amplitudes}.
They are not probabilities themselves.

\subsection{Normalization and the Born rule}

A valid state is normalized:
\[
|\alpha|^2+|\beta|^2=1.
\]
The Born rule turns amplitudes into measurement probabilities:
\[
P(0)=|\alpha|^2,\qquad P(1)=|\beta|^2.
\]
For the plus state,
\[
\ket+=\frac{\ket0+\ket1}{\sqrt2},
\]
both amplitudes are $1/\sqrt2$. Squaring their magnitudes gives
$P(0)=P(1)=1/2$. The factor $1/\sqrt2$ is therefore required; writing
$\ket0+\ket1$ without normalization would produce total probability two.

\subsection{Matrices as gates}

A matrix is a rectangular array of numbers. A one-qubit gate is represented by
a $2\times2$ complex matrix. Matrix-vector multiplication computes the new
amplitudes:
\[
\begin{bmatrix}a&b\\c&d\end{bmatrix}
\begin{bmatrix}\alpha\\\beta\end{bmatrix}
=
\begin{bmatrix}a\alpha+b\beta\\c\alpha+d\beta\end{bmatrix}.
\]
For example, the X gate exchanges the two amplitudes:
\[
\begin{bmatrix}0&1\\1&0\end{bmatrix}
\begin{bmatrix}\alpha\\\beta\end{bmatrix}
=\begin{bmatrix}\beta\\\alpha\end{bmatrix}.
\]
Applying gates in source order means multiplying by matrices in the same
temporal order: if U acts first and V second, then
\[
\ket{\psi_{\rm final}}=V(U\ket{\psi_{\rm initial}})=VU\ket{\psi_{\rm initial}}.
\]

\subsection{Requirements for a reversible gate}
\hypertarget{gate-reversibility-requirements}{}

A quantum gate is, by default, reversible: closed-system evolution is unitary,
and unitarity already implies everything ``reversible'' requires. The four
conditions below are one fact, stated four ways --- as a matrix equation, as a
statement about information, as a wire count, and as a scope limit --- so that
each can be checked on its own.

\begin{enumerate}
  \item \textbf{Unitarity.} The gate's matrix $U$ must satisfy
  \[
  U^\dagger U=UU^\dagger=I,
  \]
  where $U^\dagger$ transposes $U$ and complex-conjugates every entry
  (\emph{The complex plane, magnitude, and conjugate}), and $I$ is the
  identity matrix. Unitarity guarantees an inverse, $U^{-1}=U^\dagger$: the
  gate's action can always be undone by running $U^\dagger$. It is also
  exactly the condition that keeps a state normalized
  ($|\alpha|^2+|\beta|^2=1$, \emph{Normalization and the Born rule}) after the
  gate runs, so the total probability of all outcomes stays $100\%$.

  \item \textbf{Bijective mapping.} The gate must be a one-to-one map from
  input states to output states: given only the output and the gate, the
  input can always be reconstructed. Classical NOT has this property (an
  output of 1 proves the input was 0); classical AND and OR do not (an output
  of 0 does not say which inputs produced it), which is why they are not
  themselves reversible. A gate satisfying $U^\dagger U=I$ automatically has
  this property, because $U^\dagger$ is precisely the map that reconstructs
  the input from the output. Unitarity and bijectivity are the same fact
  viewed two ways; a reversible gate never erases information.

  \item \textbf{Equal inputs and outputs.} $U$ is a square matrix, so a
  reversible gate has exactly as many output wires as input wires: no wire
  can be discarded mid-computation without breaking the bijection above. This
  is why QPU's \code{AND} is not the two-input classical gate but a
  three-wire embedding,
  \[
  \ket{a,b,t}\mapsto\ket{a,b,t\oplus(a\land b)},
  \]
  the doubly-controlled CCNOT/Toffoli gate
  (\hyperlink{ccnot-and-embedding}{Doubly controlled NOT / Toffoli}). With
  $t$ starting at $\ket0$, $t'=a\land b$: the classical AND is recovered on
  $t'$ without erasing $a$ or $b$, at the cost of one extra wire the
  two-input classical gate did not need.

  \item \textbf{Isolation from measurement.} The three conditions above hold
  only for a closed system, with no interaction with the environment and, in
  particular, no MEASURE (\hyperlink{measure-collapse}{Measurement gate}).
  Measurement is not linear and has no inverse matrix: it collapses a
  superposition to one classical outcome and permanently discards the
  amplitudes that were not observed. RESET shares this --- preparing a wire
  at $\ket0$ regardless of its prior state is not invertible either
  (\hyperlink{reset-semantics}{Reversible gates versus RESET}). MEASURE and
  RESET are legitimate operations, just not reversible ones, and they are the
  only two exceptions to an otherwise-unitary gate list.
\end{enumerate}

